% =====================================================================
%  5.2 Arealsetningen — ekstra utfordring
% =====================================================================

\begin{infoboks}{Velg metode selv}
Arket bygger på 5.2 Arealsetningen. Vis regning og begrunnelse. Bruk eksakte svar når oppgaven er uten hjelpemidler.
Kontroller svaret ved å regne på en annen måte.
\end{infoboks}

\nivaahode{2}

\oppg{U1}{2}{\uh}
En regulær sekskant har side $4$. Del den i seks like trekanter fra midtpunktet. Finn arealet av sekskanten eksakt.

\oppg{U2}{2}{\mh}
Et trekantet beite har to sider på $50$ m og $80$ m og et areal på $1\,600$ m$^2$.
\begin{tasks}(1)
  \task Finn vinkelen mellom de to sidene. Svar med én desimal.
  \task Hvor mange ulike beiter finnes? Tegn en skisse av dem.
\end{tasks}

\oppg{U3}{2}{\uh}
Av alle trekanter der to av sidene er $7$ og $12$, hvilken har størst areal? Finn arealet og begrunn.

\nivaahode{3}

\oppg{U4}{3}{\uh}
Et parallellogram har sidene $a$ og $b$, og vinkelen mellom dem er $v$. Vis at arealet er $ab\sin v$.
Gjelder det også når $v$ er stump? Forklar.

\oppg{U5}{3}{\uh}
En regulær $n$-kant er innskrevet i en sirkel med radius $1$. Del den i $n$ trekanter fra sirkelens midtpunkt.
\begin{tasks}(1)
  \task Vis at arealet av $n$-kanten er $T_n = \dfrac{n}{2}\sin\dfrac{360^\circ}{n}$.
  \task Regn ut $T_4$, $T_6$ og $T_{12}$ eksakt.
\end{tasks}

\oppg{U6}{3}{\prog}
Se på programmet.
\begin{kode}
import math

n = 6
while n < 2000:
    T = n / 2 * math.sin(math.radians(360 / n))
    print(n, T)
    n = n * 2
\end{kode}
\begin{tasks}(1)
  \task Hvilke verdier av $T$ regner programmet ut for $n = 6$ og $n = 12$? Regn eksakt uten å kjøre det.
  \task Kjør programmet. Hvilket tall nærmer $T_n$ seg? Sammenlign med \lstinline[style=laereverk,numbers=none,frame=none]|math.pi|.
  \task Forklar hvorfor $T_n$ nærmer seg akkurat dette tallet.
  \task Hvorfor trengs \lstinline[style=laereverk,numbers=none,frame=none]|math.radians| i programmet?
  \task Endre programmet slik at det prøver $n = 3, 4, 5, \ldots$ og finner den minste $n$ der $T_n$ avviker fra $\pi$ med mindre enn $0{,}001$.
\end{tasks}

\begin{oppsummering}
Arealsetningen $T = \frac12 ab\sin C$ gjelder for alle trekanter, også med stump vinkel.
Når arealet er kjent, gir $\sin C$ to vinkler, $v$ og $180^\circ - v$, så sant sinusverdien er mindre enn $1$.
Kontroller alltid om begge vinklene gir en trekant.
\end{oppsummering}
